
\thispagestyle{empty}
\pagenumbering{gobble}

\begin{center}
    \fontsize{15pt}{1.3}\selectfont
    \textbf{ĐẠI HỌC BÁCH KHOA HÀ NỘI}
\end{center}

\vspace{3.5cm}

\begin{center}
    \fontsize{24pt}{1.3}\selectfont
    \textbf{ĐỒ ÁN II}
    
    \vspace{1cm}
    
    \fontsize{18pt}{1.5}\selectfont
    \textbf{VRAE: Cải tiến bộ mã hóa tự động}\\ \textbf{khử nhiễu và khử mờ cho biển số xe}
\end{center}

\vspace{2cm}

\begin{center}
    \fontsize{14pt}{1.5}\selectfont
    \textbf{NGUYỄN TẤT CƯỜNG}\\
    \vspace{0.5cm}
    cuong.nt227090@sis.hust.edu.vn\\
    \vspace{0.5cm}
    \textbf{Chuyên ngành: Toán Tin}
\end{center}

\vspace{3cm}

\begin{center}
    \fontsize{14pt}{1.5}\selectfont
    % Dùng 1 tabular chính. Cột cuối dùng nested tabular cho chữ ký.
    \begin{tabular}{l c l l}
        \textbf{Giảng viên hướng dẫn} & : & TS. Trần Ngọc Thăng & \begin{tabular}[t]{c} \rule{4cm}{0.5pt} \\ \footnotesize \textit{Chữ ký GVHD} \end{tabular} \\[10pt]
        \textbf{Bộ môn}               & : & Toán Tin            & \\[10pt]
        \textbf{Khoa}                 & : & Khoa Toán Tin       & \\
    \end{tabular}
\end{center}

\vfill

\begin{center}
    \fontsize{14pt}{1.3}\selectfont
    \textbf{Hà Nội, 10-2025}
\end{center}

\newpage

% --- PAGE 2: NHẬN XÉT CỦA GVHD ---
\thispagestyle{empty}

% Khung viền cho trang nhận xét (cách lề 2cm - 2.5cm)
\begin{tikzpicture}[remember picture, overlay]
    \draw[black, line width=1pt] 
        ($(current page.north west) + (2.5cm,-2.0cm)$) rectangle ($(current page.south east) + (-2.0cm,2.0cm)$);
\end{tikzpicture}

\begin{center}
    \vspace*{0.5cm} % Khoảng cách từ viền trên xuống tiêu đề
    \textbf{\fontsize{14pt}{1.3}\selectfont NHẬN XÉT CỦA GIẢNG VIÊN HƯỚNG DẪN}
\end{center}

\vspace{1cm}

\noindent \textbf{1. Mục tiêu và nội dung của đồ án}

\vspace{0.5cm}
\noindent\hspace*{1cm}(a) Mục tiêu: \dotfill

\vspace{1.5cm}
\noindent\hspace*{1cm}(b) Nội dung: \dotfill

\vspace{1.5cm}

\noindent \textbf{2. Kết quả đạt được}

\vspace{3cm}

\noindent \textbf{3. Ý thức làm việc của sinh viên}

\vfill

\begin{flushright}
    \begin{tabular}{c}
        \textit{Hà Nội, ngày \quad tháng \quad năm 2025} \\
        Giảng viên hướng dẫn \\
        \vspace{2.5cm} \\
        \textbf{TS. Trần Ngọc Thăng}
    \end{tabular} \hspace{1cm} % Căn lề phải một chút
\end{flushright}

\vspace{1cm}
\newpage

% --- PAGE 3: LỜI CẢM ƠN ---
\thispagestyle{empty}

\noindent \textbf{\fontsize{20pt}{1.3}\selectfont Lời cảm ơn}

\vspace{0.5cm}

\fontsize{15pt}{1.5}\selectfont

Em xin chân thành cảm ơn TS. Nguyễn Nam Phong và TS. Trần Ngọc Thăng đã tận tình hướng dẫn, chỉ bảo và tạo điều kiện thuận lợi cho em trong suốt quá trình thực hiện đồ án này. Những góp ý quý báu của các thầy đã giúp em hoàn thiện nghiên cứu và nâng cao chất lượng đồ án.

\vspace{0.5cm}

Em cũng xin gửi lời cảm ơn sâu sắc đến các đồng nghiệp tại Aithings đã luôn đồng hành, hỗ trợ và chia sẻ kinh nghiệm quý báu trong suốt thời gian thực hiện đồ án. Sự giúp đỡ nhiệt tình của mọi người đã tạo động lực to lớn giúp em vượt qua những khó khăn và hoàn thành tốt đồ án này.

\vspace{0.5cm}

Cuối cùng, em xin cảm ơn gia đình, bạn bè đã luôn động viên, ủng hộ và tạo điều kiện tốt nhất để em có thể hoàn thành chương trình học tập.

\vspace{2cm}

\begin{flushright}
    \begin{tabular}{c}
        \textit{Hà Nội, tháng \quad năm 2025} \\
        Tác giả đồ án \\
        \vspace{2.5cm} \\
        \textbf{Nguyễn Tất Cường}
    \end{tabular} \hspace{1cm}
\end{flushright}


\newpage

% --- PAGE 4: TÓM TẮT ĐỒ ÁN ---
\thispagestyle{empty}

\noindent \textbf{\fontsize{20pt}{1.3}\selectfont Lời nói đầu}

\vspace{1cm}

\fontsize{15pt}{1.5}\selectfont

Đồ án này là một bài nghiên cứu tập trung vào việc phát triển và thực nghiệm phương pháp \textbf{VRAE} (Vertical Residual Autoencoder) nhằm giải quyết bài toán khử nhiễu và khử mờ cho biển số xe. Đây là kết quả nghiên cứu chính trong bài báo khoa học có tiêu đề \textit{``VRAE: Vertical Residual Autoencoder for License Plate Denoising and Deblurring''} do em làm tác giả chính.

\vspace{0.5cm}

Nghiên cứu này đã được chấp nhận trình bày tại Hội nghị Quốc tế về Công nghệ Thông tin và Truyền thông lần thứ 14 \textbf{(SOICT 2025)}. Đồ án tập trung vào việc tối ưu hóa kiến trúc Autoencoder bằng cách tích hợp các kết nối dọc (Vertical Residual) để bảo toàn chi tiết vi mô và cải thiện hiệu suất khôi phục ảnh biển số xe trong các điều kiện chất lượng thấp.

\newpage
